\documentclass[11pt]{article}

\usepackage[a4paper,margin=1in]{geometry}
\usepackage{amsmath,amssymb,amsthm,mathtools}
\usepackage{enumitem}
\usepackage{microtype}

\setlength{\parskip}{0.6em}
\setlength{\parindent}{0pt}
\setlist[itemize]{leftmargin=2em}

\newcommand{\E}{\mathbb{E}}
\newcommand{\PP}{\mathbb{P}}
\newcommand{\Var}{\mathrm{Var}}
\newcommand{\1}{\mathbf{1}}
\newcommand{\dd}{\,\mathrm{d}}

\title{IE 622 \\ Detailed Solutions for Practice Problem Set 11}
\author{}
\date{}

\begin{document}

\maketitle

\section*{Problem 1}
\paragraph{Key ideas.}
This item is explicitly labeled as a reading assignment on the sheet: ``Example 3.5.4 in Norris.'' So the main purpose is to review the textbook example and absorb the method used there before solving the rest of the sheet.

\paragraph{Solution.}
Since the problem set itself does not state a theorem or a computation to be carried out, there is no standalone derivation to reproduce from the sheet alone. In a written solution set, the cleanest thing to record is that Problem 1 is background reading rather than a separate proof-based exercise.

\section*{Problem 2}
\paragraph{Key ideas.}
\begin{itemize}
    \item A stationary M/M/1 queue has geometric stationary distribution
    \[
        \pi_n = (1-\rho)\rho^n, \qquad \rho = \frac{\lambda}{\mu}, \quad n \ge 0.
    \]
    \item To study the departure count \(Y_t\), we attach a factor \(z\) every time a departure occurs and derive forward equations for the marked probabilities.
    \item The resulting generating function factorizes in a way that identifies Poisson marginals and, by induction over time intervals, independent increments.
\end{itemize}

\paragraph{Solution.}
Let \(X_t\) be the queue length and \(Y_t\) the number of departures in \([0,t]\). Assume the queue is in stationarity, so \(\rho=\lambda/\mu<1\) and \(X_0 \sim \pi\).

For \(z \in \mathbb{R}\), define
\[
    f_n(t;z) = \E_\pi\!\left[z^{Y_t}\1_{\{X_t=n\}}\right], \qquad n \ge 0.
\]
Because a departure contributes one factor of \(z\), the Kolmogorov forward equations become
\[
    \frac{\partial}{\partial t} f_0(t;z) = -\lambda f_0(t;z) + \mu z f_1(t;z),
\]
and, for \(n \ge 1\),
\[
    \frac{\partial}{\partial t} f_n(t;z)
    = \lambda f_{n-1}(t;z) + \mu z f_{n+1}(t;z) - (\lambda+\mu)f_n(t;z).
\]
The initial condition is
\[
    f_n(0;z) = \PP_\pi(X_0=n) = \pi_n.
\]

Now try the form
\[
    f_n(t;z) = \pi_n g(t;z).
\]
Substituting into the equation for \(n=0\),
\[
    \pi_0 g'(t;z)
    = -\lambda \pi_0 g(t;z) + \mu z \pi_1 g(t;z).
\]
Since \(\pi_1 = \rho \pi_0 = (\lambda/\mu)\pi_0\), this gives
\[
    g'(t;z) = \lambda(z-1)g(t;z).
\]
For \(n \ge 1\),
\[
    \pi_n g'(t;z)
    = \lambda \pi_{n-1}g(t;z) + \mu z \pi_{n+1}g(t;z) - (\lambda+\mu)\pi_n g(t;z).
\]
Using
\[
    \pi_{n-1} = \frac{1}{\rho}\pi_n = \frac{\mu}{\lambda}\pi_n,
    \qquad
    \pi_{n+1} = \rho \pi_n = \frac{\lambda}{\mu}\pi_n,
\]
the right-hand side becomes
\[
    \mu \pi_n g(t;z) + \lambda z \pi_n g(t;z) - (\lambda+\mu)\pi_n g(t;z)
    = \lambda(z-1)\pi_n g(t;z),
\]
which is the same differential equation. Since \(g(0;z)=1\), we obtain
\[
    g(t;z)=e^{\lambda(z-1)t}.
\]
Therefore
\[
    f_n(t;z)=\pi_n e^{\lambda(z-1)t}.
\]

Summing over \(n\) gives the probability generating function of \(Y_t\):
\[
    \E_\pi[z^{Y_t}] = \sum_{n \ge 0} f_n(t;z) = e^{\lambda(z-1)t}.
\]
This is exactly the pgf of a \(\mathrm{Poisson}(\lambda t)\) random variable, so
\[
    Y_t \sim \mathrm{Poisson}(\lambda t).
\]

To show that \(\{Y_t\}\) is a \emph{Poisson process}, we still need independent stationary increments. Let
\[
    0=t_0<t_1<\cdots<t_m,
    \qquad
    N_r = Y_{t_r}-Y_{t_{r-1}}.
\]
We prove by induction on \(m\) that
\[
    \E_\pi\!\left[\left(\prod_{r=1}^m z_r^{N_r}\right)\1_{\{X_{t_m}=n\}}\right]
    =
    \pi_n
    \exp\!\left(\lambda\sum_{r=1}^m (z_r-1)(t_r-t_{r-1})\right).
    \tag{*}
\]
For \(m=1\), this is exactly the identity already proved.

Assume \((*)\) holds for \(m-1\). For the last interval of length \(\Delta_m=t_m-t_{m-1}\), define
\[
    K_{i,n}(\Delta_m;z_m)
    =
    \E_i\!\left[z_m^{Y_{\Delta_m}}\1_{\{X_{\Delta_m}=n\}}\right].
\]
Using the Markov property at time \(t_{m-1}\),
\[
    \E_\pi\!\left[\left(\prod_{r=1}^m z_r^{N_r}\right)\1_{\{X_{t_m}=n\}}\right]
    =
    \E_\pi\!\left[\left(\prod_{r=1}^{m-1} z_r^{N_r}\right)K_{X_{t_{m-1}},n}(\Delta_m;z_m)\right].
\]
By the induction hypothesis with \(m-1\), the vector \((N_1,\dots,N_{m-1})\) is independent of \(X_{t_{m-1}}\), and \(X_{t_{m-1}} \sim \pi\). Hence
\[
    \E_\pi\!\left[\left(\prod_{r=1}^m z_r^{N_r}\right)\1_{\{X_{t_m}=n\}}\right]
    =
    \E_\pi\!\left[\prod_{r=1}^{m-1} z_r^{N_r}\right]
    \sum_{i \ge 0} \pi_i K_{i,n}(\Delta_m;z_m).
\]
But the sum is just the \(m=1\) formula applied over an interval of length \(\Delta_m\):
\[
    \sum_{i \ge 0} \pi_i K_{i,n}(\Delta_m;z_m)
    =
    \pi_n e^{\lambda(z_m-1)\Delta_m}.
\]
Therefore
\[
    \E_\pi\!\left[\left(\prod_{r=1}^m z_r^{N_r}\right)\1_{\{X_{t_m}=n\}}\right]
    =
    \pi_n
    \exp\!\left(\lambda\sum_{r=1}^m (z_r-1)(t_r-t_{r-1})\right),
\]
which closes the induction.

Now sum over \(n\). We get
\[
    \E_\pi\!\left[\prod_{r=1}^m z_r^{N_r}\right]
    =
    \prod_{r=1}^m \exp\!\left(\lambda(z_r-1)(t_r-t_{r-1})\right).
\]
So the increments \(N_r\) are independent, and each one is
\[
    N_r \sim \mathrm{Poisson}\bigl(\lambda(t_r-t_{r-1})\bigr).
\]
Hence \(\{Y_t\}_{t \ge 0}\) is a Poisson process of rate \(\lambda\).

\section*{Problem 3}
\paragraph{Key ideas.}
\begin{itemize}
    \item This is the \emph{immigration--death process}: arrivals occur at constant rate \(\lambda\), while each of the \(n\) active calls ends independently at rate \(\mu\), giving total completion rate \(n\mu\).
    \item For large \(t\), the effect of initial conditions dies out because each initially present call survives only with probability \(e^{-\mu t}\).
    \item For an M/M/\(\infty\) system, arrivals never wait, so a departure time is simply an arrival time shifted by an independent exponential service time. This preserves the Poisson nature of the point process of completion times.
\end{itemize}

\paragraph{Solution.}
\textbf{(a) Transition rates.}

Let \(X_t\) be the number of occupied lines. Then \(X_t\) is a CTMC on \(\{0,1,2,\dots\}\) with
\[
    q_{n,n+1} = \lambda, \qquad n \ge 0,
\]
because calls arrive as a Poisson process of rate \(\lambda\), and
\[
    q_{n,n-1} = n\mu, \qquad n \ge 1,
\]
because when \(n\) calls are active, each call ends at rate \(\mu\), so the total completion rate is \(n\mu\).

The diagonal rates are
\[
    q_{0,0}=-\lambda, \qquad q_{n,n}=-(\lambda+n\mu), \quad n \ge 1.
\]

\textbf{(b) Large-time distribution of \(X_t\).}

Fix \(t\). There are two contributions to \(X_t\):
\begin{itemize}
    \item calls already present at time \(0\) that are still active at time \(t\),
    \item calls arriving during \((0,t]\) that are still active at time \(t\).
\end{itemize}

Condition on \(X_0=n\). Each initially present call survives up to time \(t\) with probability
\[
    \PP(\text{service time}>t)=e^{-\mu t}.
\]
Therefore the number of such survivors is
\[
    S_t \mid \{X_0=n\} \sim \mathrm{Binomial}(n,e^{-\mu t}).
\]

Now consider calls arriving during \((0,t]\). An arrival at time \(s\) is still present at time \(t\) iff its service time exceeds \(t-s\), which happens with probability
\[
    e^{-\mu(t-s)}.
\]
By thinning the Poisson arrival process,
\[
    A_t \sim \mathrm{Poisson}\!\left(\int_0^t \lambda e^{-\mu(t-s)} \dd s\right)
    =
    \mathrm{Poisson}\!\left(\frac{\lambda}{\mu}(1-e^{-\mu t})\right).
\]
Moreover, \(A_t\) is independent of \(S_t\). Hence
\[
    X_t \mid \{X_0=n\}
    \overset{d}{=}
    \mathrm{Binomial}(n,e^{-\mu t})
    +
    \mathrm{Poisson}\!\left(\frac{\lambda}{\mu}(1-e^{-\mu t})\right).
\]

As \(t \to \infty\), we have \(e^{-\mu t}\to 0\), so the binomial term vanishes, while the Poisson mean tends to \(\lambda/\mu\). Therefore
\[
    X_t \Longrightarrow \mathrm{Poisson}\!\left(\frac{\lambda}{\mu}\right).
\]
So for large \(t\), \(X_t\) is approximately \(\mathrm{Poisson}(\lambda/\mu)\). In fact, under stationarity it is exactly Poisson\((\lambda/\mu)\) at every time.

\textbf{(c) Distribution of \(Y_t\).}

Under stationarity, let \(Y_t\) be the number of completed calls during \([0,t]\). Since an M/M/\(\infty\) queue has infinitely many lines, an arrival never waits. Thus if the arrival times are \(\{T_k\}\) and the service times are i.i.d.\ exponentials \(\{S_k\}\), then the departure times are
\[
    D_k = T_k + S_k.
\]
The arrivals form a Poisson point process of rate \(\lambda\), and the marks \(S_k\) are i.i.d.\ and independent of the arrival process. By the displacement/mapping theorem for marked Poisson processes, the shifted points \(\{D_k\}\) also form a Poisson point process of rate \(\lambda\).

Therefore the completion process is a Poisson process of rate \(\lambda\), and in particular
\[
    Y_t \sim \mathrm{Poisson}(\lambda t).
\]

\section*{Problem 4}
\paragraph{Key ideas.}
\begin{itemize}
    \item If a stationary distribution satisfies detailed balance, then any restriction of the chain to a subset of states inherits the same pairwise balance relations after renormalization.
    \item Once detailed balance is verified for the constrained chain, stationarity follows immediately.
\end{itemize}

\paragraph{Solution.}
Let the constrained chain on \(A\) have rates
\[
    \bar q_{x,y} =
    \begin{cases}
        q_{x,y}, & x,y \in A,\ x\ne y,\\
        0, & \text{otherwise},
    \end{cases}
\]
with diagonal terms chosen so that row sums are \(0\). Define
\[
    C = \sum_{y \in A}\pi(y),
    \qquad
    \nu(x)=\frac{\pi(x)}{C}, \qquad x \in A.
\]
Clearly \(\nu\) is a probability distribution on \(A\).

Now take any \(x,y \in A\). Because \(\pi\) satisfies detailed balance for the original chain,
\[
    \pi(x)q_{x,y} = \pi(y)q_{y,x}.
\]
Since \(\bar q_{x,y}=q_{x,y}\) for \(x,y \in A\), we have
\[
    \nu(x)\bar q_{x,y}
    =
    \frac{\pi(x)}{C} q_{x,y}
    =
    \frac{\pi(y)}{C} q_{y,x}
    =
    \nu(y)\bar q_{y,x}.
\]
So \(\nu\) also satisfies detailed balance for the constrained chain.

Now sum over \(x \in A\). For any fixed \(y \in A\),
\[
    \sum_{x \in A} \nu(x)\bar q_{x,y}
    =
    \nu(y)\sum_{x \in A}\bar q_{y,x}
    = 0,
\]
because the rows of a generator sum to \(0\). Hence \(\nu\) is stationary for \(\{Y_t\}\).

\section*{Problem 5}
\paragraph{Key ideas.}
\begin{itemize}
    \item Under stationarity, the departure process of an M/M/1 queue is Poisson with the same rate as the arrival process. This is the content of Burke's theorem, proved directly in Problem 2.
    \item Thinning a Poisson process by retaining each point independently with probability \(p\) gives another Poisson process, now of rate \(p\lambda\).
    \item The two-node system is an open Jackson network. The traffic equations determine the effective arrival rates, and the stationary law is product-form when each utilization is strictly below \(1\).
\end{itemize}

\paragraph{Solution.}
\textbf{(a) Arrival process into the second queue.}

Under stationarity, queue 1 must itself be stable, so \(\lambda<\mu_1\). By Problem 2, the departure process from the first M/M/1 queue is a Poisson process of rate \(\lambda\). Each departing customer is routed to queue 2 with probability \(p\), independently of everything else. Independent thinning of a Poisson process preserves the Poisson property, so the arrival process into queue 2 is a Poisson process of rate
\[
    p\lambda.
\]

\textbf{(b) Stability condition and stationary distribution.}

This tandem system is an open Jackson network with effective arrival rates
\[
    \beta_1 = \lambda,
    \qquad
    \beta_2 = p\lambda.
\]
Therefore the utilizations are
\[
    \rho_1 = \frac{\lambda}{\mu_1},
    \qquad
    \rho_2 = \frac{p\lambda}{\mu_2}.
\]
The queue-length process is stable if and only if both nodes are stable, namely
\[
    \lambda < \mu_1
    \qquad\text{and}\qquad
    p\lambda < \mu_2.
\]
These conditions are necessary because if either queue has input rate at least as large as its service rate, that queue cannot have a stationary queue-length distribution. They are sufficient because an open Jackson network with \(\rho_i<1\) at every node is positive recurrent.

Hence the stationary distribution is
\[
    \pi(n_1,n_2)
    =
    (1-\rho_1)(1-\rho_2)\rho_1^{n_1}\rho_2^{n_2},
    \qquad n_1,n_2 \ge 0,
\]
that is,
\[
    \pi(n_1,n_2)
    =
    \left(1-\frac{\lambda}{\mu_1}\right)
    \left(1-\frac{p\lambda}{\mu_2}\right)
    \left(\frac{\lambda}{\mu_1}\right)^{n_1}
    \left(\frac{p\lambda}{\mu_2}\right)^{n_2}.
\]

\section*{Problem 6}
\paragraph{Key ideas.}
\begin{itemize}
    \item In an open Jackson network, the \emph{traffic equations} determine the effective arrival rates \(\beta_i\) after accounting for both external arrivals and internal routing.
    \item The routing matrix has spectral radius \(<1\) in the open-network setting, so the linear system has a unique nonnegative solution obtained from a Neumann series.
    \item Jackson's theorem then gives positive recurrence and a product-form stationary distribution when every node has utilization \(\rho_i=\beta_i/\mu_i<1\).
\end{itemize}

\paragraph{Solution.}
\textbf{(a) Existence and uniqueness of the traffic equations.}

Write the traffic equations in row-vector form as
\[
    \beta = \lambda + \beta P,
\]
where \(P\) is the routing matrix. This is the same as the componentwise equation on the sheet, up to the indexing convention used for \(p_{i,j}\).

Rearranging,
\[
    \beta(I-P)=\lambda.
\]
In an open Jackson network, a customer eventually leaves the system with probability \(1\), so the routing matrix is transient in the sense that
\[
    P^n \to 0 \quad \text{as } n\to\infty.
\]
Equivalently, the spectral radius of \(P\) is \(<1\). Therefore \(I-P\) is invertible and
\[
    (I-P)^{-1} = \sum_{n=0}^{\infty} P^n.
\]
Hence the unique solution is
\[
    \beta = \lambda(I-P)^{-1}
    =
    \lambda \sum_{n=0}^{\infty} P^n.
\]
Every entry of \(P^n\) is nonnegative, so every entry of \(\beta\) is nonnegative. Thus the traffic equations have a unique nonnegative solution.

\textbf{(b) Positive recurrence and stationary distribution.}

Define
\[
    \rho_i = \frac{\beta_i}{\mu_i}, \qquad i=1,\dots,K.
\]
Assume \(\rho_i<1\) for every \(i\). Consider
\[
    \pi(n_1,\dots,n_K)
    =
    \prod_{i=1}^K (1-\rho_i)\rho_i^{n_i},
    \qquad n_i \in \{0,1,2,\dots\}.
\]
Because each \(\rho_i<1\), this is a genuine probability distribution.

Why is this the correct stationary law? The important idea is that in equilibrium each node behaves as if it were an M/M/1 queue receiving an effective arrival stream of rate \(\beta_i\). The traffic equations
\[
    \beta_i = \lambda_i + \sum_{j=1}^K \beta_j p_{j,i}
\]
say exactly that the total flow into node \(i\) equals its effective arrival rate:
\begin{itemize}
    \item external arrivals contribute \(\lambda_i\),
    \item departures from other nodes routed into node \(i\) contribute \(\sum_j \beta_j p_{j,i}\).
\end{itemize}
Under the candidate product form, the marginal law at node \(i\) is geometric with parameter \(\rho_i\), which is exactly the stationary law of an isolated M/M/1 queue with arrival rate \(\beta_i\) and service rate \(\mu_i\).

Jackson's theorem states that these local M/M/1 balances combine consistently across the whole network, yielding the global balance equations for the full CTMC. Therefore \(\pi\) is stationary.

Since \(\pi\) is a probability distribution, the invariant measure has finite total mass. For an irreducible CTMC on a countable state space, the existence of such a stationary probability distribution implies positive recurrence. Hence the network is positive recurrent, and its stationary distribution is
\[
    \pi(n)
    =
    \prod_{i=1}^K (1-\rho_i)\rho_i^{n_i},
    \qquad
    n=(n_1,\dots,n_K).
\]

\section*{Problem 7}
\paragraph{Key ideas.}
\begin{itemize}
    \item Conditional expectation given \(Z\) should be a measurable function of \(Z\), so we look for a function \(h\) with \(Y=h(Z)\).
    \item When \((X,Z)\) has joint density \(f\), the conditional density of \(X\) given \(Z=z\) is
    \[
        f_{X\mid Z}(x\mid z)=\frac{f(x,z)}{f_Z(z)},
        \qquad
        f_Z(z)=\int f(x,z)\dd x,
    \]
    whenever \(f_Z(z)>0\).
    \item The defining property of conditional expectation is verified by integrating against events in \(\sigma(Z)\), that is, sets of the form \(\{Z\in A\}\).
\end{itemize}

\paragraph{Solution.}
Define
\[
    f_Z(z)=\int f(x,z)\dd x
\]
and
\[
    h(z)=
    \begin{cases}
        \dfrac{\int g(x)f(x,z)\dd x}{f_Z(z)}, & f_Z(z)>0,\\[1.2ex]
        0, & f_Z(z)=0.
    \end{cases}
\]
Then the random variable in the problem is just
\[
    Y(\omega)=h(Z(\omega)).
\]
So \(Y\) is \(\sigma(Z)\)-measurable.

We now verify the defining identity. Let \(A\) be any Borel set. Then \(\{Z\in A\}\in \sigma(Z)\), and
\[
    \E\!\left[Y\1_{\{Z\in A\}}\right]
    =
    \int_A h(z)f_Z(z)\dd z.
\]
On the set where \(f_Z(z)>0\), \(h(z)f_Z(z)=\int g(x)f(x,z)\dd x\), and on the set where \(f_Z(z)=0\), both sides are \(0\). Hence
\[
    \E\!\left[Y\1_{\{Z\in A\}}\right]
    =
    \int_A \int g(x)f(x,z)\dd x \dd z.
\]
By Fubini's theorem, this is
\[
    \int \int g(x)\1_{\{z\in A\}} f(x,z)\dd x \dd z
    =
    \E\!\left[g(X)\1_{\{Z\in A\}}\right].
\]
Thus for every event in \(\sigma(Z)\),
\[
    \E\!\left[Y\1_B\right]=\E\!\left[g(X)\1_B\right].
\]
So \(Y\) is a version of \(\E[g(X)\mid Z]\).

\section*{Problem 8}
\paragraph{Key ideas.}
\begin{itemize}
    \item If \(B\in G\), then \(\1_B\) is \(G\)-measurable.
    \item Multiplying by a \(G\)-measurable indicator and then taking conditional expectation lets us ``localize'' identities on the set \(B\).
\end{itemize}

\paragraph{Solution.}
Let
\[
    Y_1=\E[X_1\mid G],
    \qquad
    Y_2=\E[X_2\mid G].
\]
Because \(B\in G\), the indicator \(\1_B\) is \(G\)-measurable. Therefore
\[
    \1_B Y_1 = \E[\1_B X_1 \mid G],
    \qquad
    \1_B Y_2 = \E[\1_B X_2 \mid G].
\]
But \(X_1=X_2\) on \(B\), so
\[
    \1_B X_1 = \1_B X_2 \quad \text{a.s.}
\]
Taking conditional expectation given \(G\),
\[
    \1_B Y_1 = \E[\1_B X_1\mid G] = \E[\1_B X_2\mid G] = \1_B Y_2
    \quad \text{a.s.}
\]
Hence
\[
    \E[X_1\mid G] = \E[X_2\mid G]
    \quad \text{a.s. on } B.
\]

\section*{Problem 9}
\paragraph{Key ideas.}
\begin{itemize}
    \item The easiest way to prove identities for conditional expectation is to verify the defining properties:
    \begin{itemize}
        \item \(G\)-measurability,
        \item equality of integrals over all sets \(B\in G\).
    \end{itemize}
    \item Positivity ultimately comes from the fact that the expectation of a nonnegative random variable cannot be negative on any \(G\)-measurable event.
\end{itemize}

\paragraph{Solution.}
\textbf{(a) Linearity.}

Let
\[
    Y=\alpha \E[X_1\mid G] + \beta \E[X_2\mid G].
\]
Since conditional expectations are \(G\)-measurable, \(Y\) is also \(G\)-measurable.

Now fix any \(B\in G\). Then
\[
    \E[Y\1_B]
    =
    \alpha \E\!\left[\E[X_1\mid G]\1_B\right]
    +
    \beta \E\!\left[\E[X_2\mid G]\1_B\right].
\]
By the defining property of conditional expectation,
\[
    \E\!\left[\E[X_1\mid G]\1_B\right]=\E[X_1\1_B],
    \qquad
    \E\!\left[\E[X_2\mid G]\1_B\right]=\E[X_2\1_B].
\]
So
\[
    \E[Y\1_B]
    =
    \alpha \E[X_1\1_B]+\beta \E[X_2\1_B]
    =
    \E[(\alpha X_1+\beta X_2)\1_B].
\]
Therefore \(Y\) satisfies the defining properties of \(\E[\alpha X_1+\beta X_2\mid G]\), and hence
\[
    \E[\alpha X_1+\beta X_2\mid G]
    =
    \alpha \E[X_1\mid G]+\beta \E[X_2\mid G].
\]

\textbf{(b) Positivity.}

Suppose \(X_1 \le X_2\) almost surely. Then
\[
    X:=X_2-X_1 \ge 0 \quad \text{a.s.}
\]
We first claim that
\[
    \E[X\mid G] \ge 0 \quad \text{a.s.}
\]
Indeed, if the set
\[
    B=\{\E[X\mid G]<0\}
\]
had positive probability, then \(B\in G\) and
\[
    \E[X\1_B]
    =
    \E[\E[X\mid G]\1_B]
    <0,
\]
which is impossible because \(X\1_B\ge 0\) almost surely. So \(\E[X\mid G]\ge 0\) a.s.

Now use linearity:
\[
    \E[X_2\mid G]-\E[X_1\mid G]
    =
    \E[X_2-X_1\mid G]
    =
    \E[X\mid G]
    \ge 0
    \quad \text{a.s.}
\]
Hence
\[
    \E[X_1\mid G] \le \E[X_2\mid G]
    \quad \text{a.s.}
\]

\section*{Problem 10}
\paragraph{Key ideas.}
\begin{itemize}
    \item Condition on \(N\). Given \(N=n\), the random sum becomes an ordinary sum of \(n\) i.i.d.\ variables.
    \item Use the law of total expectation:
    \[
        \E[X]=\E[\E[X\mid N]].
    \]
    \item Use the law of total variance:
    \[
        \Var(X)=\E[\Var(X\mid N)] + \Var(\E[X\mid N]).
    \]
\end{itemize}

\paragraph{Solution.}
Let
\[
    X = Y_1+\cdots+Y_N.
\]
Because \(N\) is independent of the i.i.d.\ sequence \(\{Y_i\}\), once we condition on \(N=n\), the random variable \(X\) is the sum of \(n\) i.i.d.\ copies of \(Y_1\).

\textbf{(a) Mean.}

Given \(N=n\),
\[
    \E[X\mid N=n] = n\mu.
\]
Therefore
\[
    \E[X\mid N]=\mu N.
\]
Now apply total expectation:
\[
    \E[X] = \E[\E[X\mid N]] = \E[\mu N] = \mu \E[N].
\]

\textbf{(b) Variance.}

Given \(N=n\),
\[
    \Var(X\mid N=n)=n\sigma^2,
\]
since variances of independent summands add. Hence
\[
    \Var(X\mid N)=N\sigma^2.
\]
Also, from part (a),
\[
    \E[X\mid N]=\mu N.
\]
Now apply the law of total variance:
\[
    \Var(X)
    =
    \E[\Var(X\mid N)] + \Var(\E[X\mid N]).
\]
Substituting the two expressions above,
\[
    \Var(X)
    =
    \E[N\sigma^2] + \Var(\mu N)
    =
    \sigma^2 \E[N] + \mu^2 \Var(N).
\]
That is exactly the required identity.

\end{document}
